% Cuestión III. Procedimiento de estegoanálisis de los 4 archivos del grupo 5.
% Cifras de generado/datos-analisis.tex (informe/scripts/medir_estegoanalisis.py).

\subsection{Procedimiento}

Los cuatro archivos (\texttt{avatar.bmp}, \texttt{madrid.bmp}, \texttt{montevideo.bmp} y \texttt{secreto1.bmp}) se analizaron con los mismos pasos:
\begin{enumerate}
\item \textbf{Cabecera y tamaño.} Se comparó el tamaño real del archivo con el que corresponde a la cabecera más los píxeles. Sólo \texttt{madrid.bmp}
tiene bytes de más (\dato{an-madrid-agregados}).
\item \textbf{Prueba de los tres métodos.} Se ejecutó \texttt{stegobmp -extract} con LSB1, LSB4 y LSBI sin contraseña. Con el método equivocado el campo de
tamaño (4 bytes big endian) sale como un número al azar, casi siempre mayor que la capacidad; con el correcto cabe en la capacidad y los datos
terminan en una extensión válida (cuadro~\ref{tab:q3-metodos}).
\item \textbf{Examen de lo extraído.} Se identificó el tipo por su firma (\texttt{tools/sigsearch.py}) y se revisó si había datos después de la marca de fin
del formato.
\item \textbf{Cifrado.} Si el tamaño es válido pero no hay extensión, el contenido está cifrado; los parámetros salen de las pistas de los otros archivos.
\end{enumerate}

\begin{table}[htbp]
\centering
\small
\renewcommand{\arraystretch}{1.1}
\begin{tabularx}{\textwidth}{@{}>{\raggedright\arraybackslash}p{2.6cm}>{\raggedright\arraybackslash}X>{\raggedright\arraybackslash}X>{\raggedright\arraybackslash}X@{}}
\toprule
Archivo & LSB1 & LSB4 & LSBI \\
\midrule
\texttt{avatar.bmp} & tamaño 1\,610\,612\,738 no cabe & tamaño 17\,825\,792 no cabe & \textbf{PDF}, 20\,322 bytes \\
\texttt{madrid.bmp} & tamaño 1\,427\,081\,616 no cabe & tamaño 1\,877\,846\,507 no cabe & tamaño 3\,275\,516\,423 no cabe \\
\texttt{montevideo.bmp} & \textbf{PNG}, 42\,925 bytes & tamaño 1\,814\,619\,278 no cabe & tamaño 14\,460\,194 no cabe \\
\texttt{secreto1.bmp} & tamaño 501\,759\,586 no cabe & tamaño válido, extensión malformada & tamaño 3\,161\,987\,507 no cabe \\
\bottomrule
\end{tabularx}

\caption{Resultado de \texttt{stegobmp -extract} con cada método, sin contraseña}
\label{tab:q3-metodos}
\end{table}

\subsection{Qué se encontró}

\begin{table}[htbp]
\centering
\small
\renewcommand{\arraystretch}{1.15}
\begin{tabularx}{\textwidth}{@{}>{\raggedright\arraybackslash}p{2.6cm}>{\raggedright\arraybackslash}p{3.6cm}>{\raggedright\arraybackslash}X@{}}
\toprule
Archivo & Método & Contenido \\
\midrule
\texttt{avatar.bmp} & LSBI, sin cifrar & PDF de \dato{an-avatar-pdf-bytes} bytes: <<La password es metadata>>. \\
\texttt{montevideo.bmp} & LSB1, sin cifrar & PNG de \dato{an-mon-png-bytes} bytes con un tablero de buscaminas y un ZIP concatenado (Cuestión~IV). \\
\texttt{secreto1.bmp} & LSB4, \texttt{aes256} \texttt{ofb} & Video \texttt{.wmv} de \dato{an-secreto1-wmv-bytes} bytes (Cuestión~V). \\
\texttt{madrid.bmp} & Texto agregado al final del archivo & Instrucción para convertir el PNG en ZIP (Cuestión~VI). \\
\bottomrule
\end{tabularx}
\caption{Contenido recuperado de cada archivo}
\label{tab:q3-contenido}
\end{table}

En \texttt{secreto1.bmp} el tamaño de LSB4 (\dato{an-secreto1-tam4}) es válido pero no hay extensión, y no es múltiplo de 8 ni de 16: el contenido está
cifrado con un modo de flujo (CFB u OFB). En \texttt{madrid.bmp} ningún método da un tamaño coherente; el mensaje está en los bytes agregados al final.
